\subsection{Small act cones by quotient modules}

\begin{lemma}[Proved: at most two outside elements]
\label{CP2:A-two-outside-acts}
Let $M$ be a finite monoid with minimal ideal $K$ and
$|M\setminus K|\leq2$. For every finite-dimensional right
$R=\mathbb QM$-module $L$, eventual vanishing of
$\operatorname{Tor}^R_j(L,\mathbb Q)$ implies its vanishing
for every $j\geq3$. Thus a finite act cone over $M$ with
bounded rational homology has zero fourth rational homology.
\end{lemma}
\begin{proof}
If $K=M$, the monoid is a finite group and Maschke's
averaging makes its rational algebra semisimple. Otherwise
put $A=R/\mathbb QK$. At a minimal idempotent $e$, with
maximal group $G$, the ideal factorization is
$\mathbb QK=Re\otimes_{\mathbb QG}eR$.
The second factor is a free left $\mathbb QG$-module,
so this is projective as a left $R$-module. Hence
$\operatorname{pd}_R A\leq1$.
Averaging $G$ gives a projective row module $P$ and
$0\to V\to P\to\mathbb Q\to0$, with $KV=0$.
For $j\geq2$, dimension shifting identifies
$\operatorname{Tor}_j^R(L,\mathbb Q)$ with
$\operatorname{Tor}_{j-1}^R(L,V)$.

A unital rational algebra of dimension at most two is
semisimple or the dual numbers. In the first case $V$
is an $A$-projective module and has $R$-projective
dimension at most one. In the second case
$V=A^a\oplus k^b$, where $k=A/(\varepsilon)$.
The exact sequence $0\to k\to A\to k\to0$ and
$\operatorname{pd}_R A\leq1$ give
\[
 \operatorname{Tor}_j^R(L,k)
   \cong\operatorname{Tor}_{j-1}^R(L,k),\qquad j\geq3.
\]
If $b=0$ the assertion is immediate; otherwise the
eventually zero Tor for $V$ makes the displayed groups
eventually zero, hence zero already in degree two.
Shifting back proves the assertion. For an act cone
use $L=\ker(\mathbb QS\to\mathbb Q)$ and its reduced
Tor formula, with the one-degree suspension.
\end{proof}

\begin{lemma}[Proved: a connected graph with one cycle]
\label{CP2:A-one-cycle-graph-cofiber}
Let a finite monoid $M$ act cellularly on a connected
finite graph $\Gamma$, with $H_1(\Gamma;\mathbb Q)=\mathbb Q$.
For $X=\operatorname{hocofib}(\Gamma_{hM}\to BM)$,
\[
 \widetilde H_n(X;\mathbb Q)
  =\operatorname{Tor}_{n-2}^{\mathbb QM}
           (H_1(\Gamma;\mathbb Q),\mathbb Q),\qquad n\geq2.
\]
Nonzero $H_2(X;\mathbb Q)$ forces every actor to have
degree one on the graph cycle. If a commuting graph
self-map has degree $d\ne0,1$, all positive rational
cohomology products on $X$ vanish.
\end{lemma}
\begin{proof}
The augmented graph complex, with its actual right module
structure, has only its first homology as nonzero reduced
homology. Inclusion of that kernel is an equivariant
quasi-isomorphism. Tensoring a free left monoid resolution
and taking the cofiber shifts this kernel by two, proving
the formula. Its degree-two group is the coinvariant space
of the one-dimensional cycle module. Any actor degree
different from one annihilates that coinvariant space.
A commuting degree-$d$ map acts by $d$ on this kernel,
hence on all positive cofiber homology and cohomology.
For positive classes $u,v$, naturality gives $d(uv)=d^2uv$,
which proves the last assertion over $\mathbb Q$.
\end{proof}

\begin{proposition}[Proved: four states opposite a right-zero corner]
\label{CP2:A-right-zero-four-states}
Let $R_3=\{1,p,q\}$ with $ab=b$ for $a,b\in\{p,q\}$.
If a left-$R_3$, right-$M$ biset with at most four states
has collage of type $\CP^2$, an ordinary right-$M$ act
cone on at most four states also has that type.
Consequently no such collage exists when $|M|\leq12$.
For any three-element left corner, the profile
$(|M|,3,4)=(9,3,4)$ is excluded as well.
\end{proposition}
\begin{proof}
One or two states are excluded by the singleton join and
the universal two-cross theorem. Three states use
Theorem~\ref{CP2:A-three-cross-reduction}. Suppose there
are four. The two left actors $P,Q$ are retractions onto
the same nonempty subset $U$, by $PQ=Q$ and $QP=P$.
The natural graph model of the left action category has
vertices $U$ and one edge for each $s\notin U$, directed
from $Ps$ to $Qs$. Its construction and equivariance were
proved in Lemma~\ref{CP2:A-high-P5-graph}. A commuting
right actor maps an edge at its unit parameter to the
edge of the image state, or collapses it to a vertex.
Since $BR_3$ is contractible, the collage is the cofiber
of this graph's homotopy-orbit map to $BM$.

If $|U|=1$, its unique point is fixed by both sides,
and the relative-cup obstruction applies. If $P=Q$,
including $|U|=4$, the common map is a biset retraction
onto $U$. The retract is equivalent to the ordinary
$M$-act cone on $U$: the left side acts trivially and
its contractible classifying space can be collapsed
in the product edge bar and at its vertex. By
Lemma~\ref{CP2:A-biset-retract} this retract retains
the projective-plane type.

If $|U|=3$ and $P\ne Q$, the single graph edge joins
two distinct vertices. Its two components are
contractible. The equivariant map to its component
set is a homotopy equivalence, including after taking
the monoid bar. Its cofiber is therefore an ordinary
two-state act cone, already excluded.

The remaining graph has two vertices and two edges.
If it is disconnected, both edges are loops and
$P=Q$, already treated. If connected it has one
cycle. A putative projective plane forces every
actor to have degree one by the preceding lemma.
For one loop and one connecting edge, this forces
the unique loop to map to itself; its vertex is
a common fixed state. For two connecting edges
with the same direction, a degree-one actor fixes
both vertices and edges: swapping the edges has
degree minus one and folding or collapsing gives
degree zero. Again there is a common fixed state.
If the two edges have opposite directions, the
only degree-one actors are the identity and the
half-turn. Identify this graph with the circle
so that the half-turn is $z\mapsto-z$. The map
$z\mapsto z^3$ commutes with both actors. The
preceding degree lemma kills all positive products.
This exhausts the four-state graph types.

The resulting ordinary cones on one or two states
have zero products; on three states use the exact
minimum twenty-two; on four states use
Corollary~\ref{CP2:A-four-state-small}. This proves
the assertion for $|M|\leq12$.
Finally, every three-element monoid other than
$R_3$ has projective trivial right rational module.
Indeed its minimal ideal is a group, has one column,
or is the two-column one-row band with an identity;
the last case is exactly $R_3$. In the first two
cases finite-group averaging supplies the splitting
idempotent. Exact orbit comparison thus replaces
the left corner of $(9,3,4)$ by a trivial corner,
preserving rational homology and the cohomology ring.
The bounded-nine theorem excludes that orbit act
cone. This proves the last assertion without a
new multiplication census.
\end{proof}

\begin{lemma}[Proved: the two-column quotient has projective kernel]
\label{CP2:A-two-column-orbit-kernel}
Let the minimal ideal of $M$ be the band $I\times\{0,1\}$,
and suppose a finite nonempty right act $S$ has no point
fixed by the whole band. Let $o$ be the number of core
orbits in $SK$, and let $\bar S$ be the constant-action
quotient of Lemma~\ref{CP2:A-constant-quotient-H2}.
Over $R=\mathbb QM$ there is an exact sequence
\[
 0\to(eR)^{o-1}\to L_S\to L_{\bar S}\to0,
 \qquad L_T=\ker(\mathbb QT\to\mathbb Q).
\]
Every core orbit has two states, and
$|\bar S|=|S|-2(o-1)$.
Writing $X,Y$ for the two act cones, their actual quotient
map induces isomorphisms in rational homology in degrees
at least three and an exact sequence
\[
 0\to H_2(X)\to H_2(Y)\to\mathbb Q^{o-1}
       \to H_1(X)\to H_1(Y)\to0.
\]
If $X$ has the rational homology and ring of $\CP^2$,
then $H_2(Y)=\mathbb Q^o$, $H_4(Y)=\mathbb Q$, and
all its other positive rational homology vanishes.
Its degree-two cup pairing has rank exactly one.
\end{lemma}
\begin{proof}
A core orbit is a quotient of the two-point column act.
It cannot be a singleton by the hypothesis, so label it
$\{s_{a,0},s_{a,1}\}$, $0\leq a<o$. Every outside actor
has the same column map on every such orbit:
$s_{a,j}m=s_{a,\tau_m(j)}$. In particular the orbits
are individually $M$-invariant.

The kernel $L_SK$ of the quotient map has basis
$s_{a,j}-s_{0,j}$ for $1\leq a<o$, $j=0,1$.
Each is a generating relation after multiplication
by a band element in column $j$. Conversely every
generating relation is a difference of column-$j$
points, so this is the whole kernel. Its action formula
identifies each pair with $eR$, the projective two-point
row ideal. Noncore states are never identified, giving
the cardinality formula.

Apply the Tor exact sequence and
$eR\otimes_R\mathbb Q=\mathbb Q$ to the reduced
act-cone chain formula. This proves both the higher
isomorphisms and the stated low-degree sequence.
For a rational projective plane, $H_1(X)=0$ gives
the asserted dimensions, and the higher groups
transfer unchanged. Pullback is surjective in degree
two and an isomorphism in degree four. Naturality
therefore identifies the cup form on $Y$ with the
pullback of the nonzero rank-one form on $X$.
Finally compression gives $Y\simeq BN$ with minimal
band $I\times\bar S$ and the same outside elements
as $M$, so the band--Tor formula applies to this
actual quotient without any separate boundedness
assumption on $BM$.
\end{proof}

\begin{lemma}[Proved: a two-dimensional module over a small local algebra]
\label{CP2:A-local-three-small-module}
Let $A$ be a local rational algebra with $\dim A\leq3$.
If one of two finite modules $W,V$ has dimension at most
two, eventual vanishing of $\operatorname{Tor}^A_j(W,V)$
implies $\operatorname{Tor}^A_2(W,V)=0$.
\end{lemma}
\begin{proof}
By opposition suppose $\dim V\leq2$. If $V$ is
semisimple, it is a sum of copies of the residue
division algebra. A nonprojective $W$ then has
nonzero Tor in every positive degree by
Lemma~\ref{A:local-target-facts}; a projective $W$
has none. If $V$ is not semisimple, its dimension
is exactly two, its residue division algebra is
$\mathbb Q$, and its top has dimension one.
Thus it is cyclic. When $\dim A=2$ its cyclic
cover is an isomorphism. When $\dim A=3$ the
kernel of that cover is the one-dimensional simple
module $k$. For $j\geq2$, dimension shifting gives
$\operatorname{Tor}_j(W,V)=\operatorname{Tor}_{j-1}(W,k)$.
The same local-algebra fact excludes eventual
vanishing with a nonzero degree-two group. These
cases exhaust the possible dimensions.
\end{proof}

\begin{lemma}[Proved: the local rank-two additive obstruction]
\label{CP2:A-local-four-rank-two}
Let $A$ be local of dimension four, with
$\dim V=2$ and $\dim W\leq3$. There is no module
pair with bounded Tor, tensor dimension two,
$\operatorname{Tor}_1(W,V)=0$ and
$\operatorname{Tor}_2(W,V)\ne0$.
\end{lemma}
\begin{proof}
If $V$ is semisimple, boundedness and
Lemma~\ref{A:local-target-facts} force $W$ projective,
contrary to nonzero positive Tor. Otherwise the residue
field is $\mathbb Q$ and $V$ is cyclic of dimension
two. Its cyclic projective cover has nonzero kernel
$J_V$ of dimension two. The Tor exact sequence gives
\[
 0\to W\otimes_AJ_V\to W\to W\otimes_AV\to0,
 \qquad \dim(W\otimes_AJ_V)=\dim W-2.
\]
The tensor product of two nonzero modules over a local
finite algebra is nonzero, by its quotient onto the
tensor of their tops. Hence $\dim W=3$ and the
displayed tensor dimension is one. The same top
argument makes $W$ cyclic. Its cyclic cover then
has one-dimensional simple kernel. Consequently
$\operatorname{Tor}_j(W,V)=\operatorname{Tor}_{j-1}(k,V)$
for $j\geq2$, which is nonzero in arbitrarily high
degrees because $V$ is not projective. This contradicts
boundedness.
\end{proof}

\begin{theorem}[Proved: bounded act cones through monoid order nine]
\label{CP2:A-bounded-nine-acts}
For every finite monoid of order at most nine and every
nonempty finite right act, bounded rational homology
of its act cone forces every rational degree-two
product to vanish. There is no bound on the number
of states or on the second Betti number. The left-act
version holds by opposition.
\end{theorem}
\begin{proof}
If $H^2(BM;\mathbb Q)=0$, use the suspension quotient.
If there are at most two minimal rows, use
Lemma~\ref{CP2:A-high-two-row-act}. In the remaining
cases the minimal ideal has at least three rows and
two columns. Cardinality at most nine forces trivial
maximal group and either a $3\times2$ band with
at most three outside elements or a $4\times2$
band with one outside element. Two or fewer outside
elements are covered by Lemma~\ref{CP2:A-two-outside-acts}.

Only the $3\times2$ band with three outside elements
remains. A point fixed by all of $K$ is fixed by $M$:
if $ze=z$, then $zm=z(em)=z$, since $em\in K$.
It gives a section of the action bar and makes
all second classes relative, so their products vanish.
Otherwise the orbit-kernel lemma and compression give
a monoid $N$ with bounded rational homology and
$H_4(BN;\mathbb Q)=H_4(X;\mathbb Q)$. Its contracted
algebra $A$ has dimension three and its row
augmentation module has dimension two. For nonlocal
$A$, Lemma~\ref{A:nonlocal-four} forces degree-two
Tor to vanish under boundedness. For local $A$ use
Lemma~\ref{CP2:A-local-three-small-module}. The
band--Tor formula therefore gives $H_4(X;\mathbb Q)=0$.
This also annihilates every degree-two product.
\end{proof}

\begin{proposition}[Proved: ten-element corners and at most six states]
\label{CP2:A-ten-six-state-rational}
An act cone over a monoid of order at most ten, with
at most six states, cannot have the rational homology
and cohomology ring of $\CP^2$.
\end{proposition}
\begin{proof}
The preceding theorem handles orders through nine.
The relative and two-row obstructions leave, at order
ten, band shapes $3\times2$, $4\times2$, $3\times3$.
The last two have at most two outside elements and
are excluded by Lemma~\ref{CP2:A-two-outside-acts}.
Thus $K=3\times2$ and there are four outside elements.
A band-fixed state is a global fixed state and gives
zero products as above. Otherwise let $o$ count the
two-point core orbits. The act must be connected:
the cofiber exact sequence identifies $H_1(X;\mathbb Q)$
with the reduced zeroth homology of its action category,
since $H_1(BM;\mathbb Q)=0$. Because each core orbit
is invariant, $o=3$ at six or fewer states would have
$S=SK$ and a disconnected act. Hence $o=1$ or $2$.

The quotient compression has a four-dimensional
contracted algebra $A$, row module $V$ of dimension
two, and column module $W$ of dimension
$|S|-2(o-1)-1$. Its required Tor dimensions are
$(o,0,1,0,\ldots)$. Nonlocal $A$ is excluded by
Lemma~\ref{A:nonlocal-four}. If $A$ is local and
$o=1$, then $\dim W\leq5$, and
Lemma~\ref{A:local-four-small-target} excludes the
required four initial Tor dimensions. If $o=2$,
then $\dim W\leq3$, and
Lemma~\ref{CP2:A-local-four-rank-two} applies.
This excludes the stated rational target in all cases.
\end{proof}

\begin{lemma}[Proved: four-state band acts with zero first homology]
\label{CP2:A-four-state-band}
Let $M$ have a rectangular-band minimal ideal, and
let $S$ be a right act with at most four states.
If $H_1(B\mathcal C(M,S);\mathbb Q)=0$, either all
rational degree-two products vanish or every band
actor is constant on $S$.
\end{lemma}
\begin{proof}
A band-fixed point is globally fixed and gives zero
products. Otherwise every core orbit has at least
two points. Every such orbit is individually invariant
under $M$, since $yK$ is invariant and contains $y$
for a core point. Two or more core orbits at four
states force exactly two two-point orbits and $S=SK$.
The action category would be disconnected, contrary
to the cofiber first-homology sequence. Thus there
is just one core orbit.

On any right orbit of a rectangular band, each band
element acts constantly: writing a point as $yk$,
the value $(yk)\ell$ depends only on the fixed row
of the orbit and the column of $\ell$. For arbitrary
$s\in S$, idempotence gives $s\ell=(s\ell)\ell$,
and $s\ell$ lies in this single core orbit. Hence
the constant value is independent of $s$ as well.
\end{proof}

\begin{corollary}[Proved: four-state act cones through seventeen arrows]
\label{CP2:A-four-state-small}
An ordinary four-state act cone over a monoid of order
at most twelve cannot have type $\CP^2$.
\end{corollary}
\begin{proof}
Nonzero square requires nonzero corner $H^2$, at least
two minimal columns and at least three minimal rows,
by the relative and two-row results. A nontrivial
minimal group would cost at least $3\cdot2\cdot2+1=13$
elements. Thus the ideal is a band. The preceding
lemma forces constant action. Compression gives
a monoid of order
\[
 m-rs+4r\leq m+2r
 \leq m+2\lfloor(m-1)/2\rfloor\leq22
 \qquad(m\leq12).
\]
This contradicts the complete monoid minimum.
\end{proof}
